\chapter{Cadre et contexte du projet}
\label{chap:contexte}

Dans ce chapitre, nous présentons le contexte dans lequel s'inscrit ce
projet, la problématique qui en découle, les objectifs poursuivis,
l'intérêt de l'étude ainsi que ses limites.

\section{Contexte}

Les réseaux mobiles de quatrième et cinquième génération (4G/5G)
reposent sur une architecture qui sépare le réseau d'accès radio
(\emph{Radio Access Network}, RAN), constitué des stations de base
desservant les terminaux des utilisateurs, et le c\oe ur de réseau
(\emph{Core Network}), qui gère l'authentification, la mobilité et
l'acheminement du trafic vers les services (figure~\ref{fig:5g-arch}).
Chaque station de base couvre une zone géographique limitée et rapporte
en continu des indicateurs de trafic (volume de données, nombre
d'appels, de SMS) qui peuvent être agrégés par zone et par créneau
horaire.

\begin{figure}[htbp]
  \centering
  \includegraphics[width=0.7\textwidth]{5g-architecture.png}
  \caption[Architecture simplifiée d'un réseau 5G]{Architecture simplifiée d'un réseau 5G : terminaux (UE), réseau d'accès radio (5G NR) et c\oe ur de réseau (5G Core). Source~: GiftedYeti, \emph{5G Architecture.png}, Wikimedia Commons, licence CC BY-SA 4.0.}
  \label{fig:5g-arch}
\end{figure}

Ce volume massif de données d'usage, collecté en continu sur
l'ensemble d'une métropole, constitue une matière première précieuse
pour la supervision du réseau. Il permet en principe de repérer, à
l'échelle d'une zone géographique, des journées où le trafic s'écarte
significativement de son comportement habituel -- signe possible d'un
incident technique, d'une saturation, ou au contraire d'un événement
extérieur (jour férié, rassemblement, manifestation sportive) qui
modifie temporairement les habitudes de déplacement et de
communication des usagers.

\section{Problématique}

La difficulté principale réside dans la distinction entre une
variation normale du trafic -- heures de pointe, week-ends,
saisonnalité -- et un comportement réellement atypique. À l'échelle
d'une métropole, cette analyse doit en outre être menée sur un grand
nombre de zones géographiques et sur plusieurs canaux de communication
(Internet, SMS, appels), dont les dynamiques diffèrent sensiblement
les unes des autres. Une méthode de détection doit donc être
suffisamment sensible pour repérer les journées atypiques, tout en
restant robuste au bruit statistique inhérent à des données réelles à
large échelle, et sans disposer d'anomalies déjà étiquetées pour
entraîner un modèle supervisé.

\section{Objectifs}

\subsection{Objectif général}

L'objectif général de ce travail est de concevoir et d'évaluer une
méthode de détection d'anomalies dans le trafic de réseaux cellulaires
4G/5G, appliquée à des données réelles de trafic mobile de la ville de
Milan.

\subsection{Objectifs spécifiques}

Plus spécifiquement, ce travail vise à~:
\begin{itemize}
  \item encoder le profil horaire de trafic de chaque zone sous une
        forme compacte et comparable d'un jour à l'autre~;
  \item détecter, pour chaque zone, les journées présentant un
        comportement statistiquement atypique~;
  \item appliquer cette détection indépendamment sur les trois canaux
        disponibles (Internet, SMS, appels) et évaluer l'apport d'une
        confirmation multi-canal~;
  \item valider les anomalies détectées par recoupement avec un
        calendrier d'événements connus (jours fériés, matchs à
        domicile à San Siro).
\end{itemize}

\section{Intérêt de l'étude}

\subsection{Intérêt académique}

Sur le plan académique, ce travail permet d'approfondir des notions de
traitement de séries temporelles, d'algorithmique des chaînes de
caractères (distance d'édition) et d'analyse statistique, en les
confrontant à un cas d'usage concret et à des données réelles de
grande échelle. Il s'inscrit également dans la continuité des travaux
de recherche existants sur la détection d'anomalies dans les réseaux
cellulaires, dont le framework DiNATrAX constitue un point de départ
méthodologique.

\subsection{Intérêt professionnel}

Sur le plan professionnel, une méthode de détection d'anomalies
fiable, applicable à plusieurs canaux de trafic et ne nécessitant pas
de données étiquetées, présente un intérêt direct pour la supervision
opérationnelle des réseaux mobiles~: elle peut aider à prioriser les
vérifications terrain, à anticiper des besoins de capacité localisés,
ou à documenter automatiquement des journées atypiques pour les
équipes d'exploitation.

\section{Limites de l'étude}

Cette étude présente plusieurs limites. Premièrement, les données
utilisées portent sur une période de deux mois (novembre-décembre
2013) et sur une seule ville, ce qui limite la généralisation directe
des seuils calibrés à d'autres contextes géographiques ou temporels.
Deuxièmement, la validation des anomalies détectées repose sur un
recoupement avec un calendrier d'événements connus, qui ne permet de
confirmer qu'une partie des anomalies observées -- les anomalies sans
explication calendaire identifiée ne sont pas nécessairement des faux
positifs, mais restent plus difficiles à interpréter. Enfin, le jeu de
données ne distingue pas les usages 4G des usages 5G à proprement
parler, la technologie 5G n'étant pas encore déployée sur la période
couverte~; l'approche reste néanmoins directement transposable à des
données de trafic 5G actuelles.
